\input{../tex-inputs/latex-leseansicht-vorspann}

\section[Gustav Schwarzkopf an Arthur Schnitzler, 25. 11. 1894]{L04279 Gustav Schwarzkopf an Arthur Schnitzler, 25. 11. 1894}
\nopagebreak\mylabel{L04279v}
\rehead{ }\normalsize\beginnumbering\briefempfaengerindex{Schnitzler, Arthur@\textsc{Schnitzler, Arthur}!zzzSchwarzkopf, Gustav@\emph{von Gustav Schwarzkopf}!1894-11-252@{25. 11. 1894}|(be}
\toendnotes[C]{\smallbreak\pagebreak[2]}
\correspDesc{Versand  durch Gustav Schwarzkopf am 25. 11. 1894 in Wien
\newline{}Erhalt  durch Arthur Schnitzler im Zeitraum [25. 11. 1894 – 28. 11. 1894?] in Wien}\toendnotes[C]{\smallbreak}
\Standort{CUL, Schnitzler, B 96a.}
\physDesc{Briefkarte, 277 Zeichen (Briefkarte mit seitlichem Papierabriss)
\newline{}Handschrift: schwarze Tinte, deutsche Kurrent}\toendnotes[C]{\smallbreak}
\pstart
           \raggedleft{}{\pb}25. XI 94\pend
           
\pstart{}Lieber Freund!\pend\vspace{0.5em}
\pstart
           herzlichen Dank für ihr Buch\pwindex{Schnitzler, Arthur 15.\,5.\,1862 Wien – 21.\,10.\,1931 ebd.@\textsc{Schnitzler, Arthur} (15.\,5.\,1862 Wien – 21.\,10.\,1931 ebd.), \emph{Schriftsteller, Mediziner}!Sterben. Novelle@\strich\emph{Sterben. Novelle}|pwv}.
      Es iſt eine{ }ſchöne wertvolle
      Arbeit, die in ihrer vornehmen,
      echt künſtleriſchen Einfachheit
      ergreifend wirkt. Ich finde
            {\pb}hoffentlich wol bald Gelegenheit
      mit Ihnen darüber zu{ }ſprechen.\pend
           
\pstart
           Herzlichſt{\\[\baselineskip]}Ihr{\\[\baselineskip]}aufrichtig ergebener{\\[\baselineskip]}\spacefill\mbox{GustavSch}\pend
           \leftskip=0em{}\selectlanguage{ngerman}\endnumbering\briefempfaengerindex{Schnitzler, Arthur@\textsc{Schnitzler, Arthur}!zzzSchwarzkopf, Gustav@\emph{von Gustav Schwarzkopf}!1894-11-252@{25. 11. 1894}|)be}\mylabel{L04279h}
\begin{anhang}
\end{anhang}\newcommand{\dateiname}{L04279}\newcommand{\titel}{Gustav Schwarzkopf an Arthur Schnitzler, 25. 11. 1894}\newcommand{\editorInnen}{Herausgegeben von Jahnke, SelmaMüller, Martin Anton}\input{../tex-inputs/latex-leseansicht-abspann}
